\section{Doctrine of Castes}

\label{sec:DoctrineCastes}

\begin{quotex}
One law for the lion and ox is oppression.
\flright{\textsc{William Blake}, \emph{Marriage of Heaven and Hell}}
\end{quotex}


The characteristics of each caste is shown in chart \ref{tab:CharacteristicsCaste}, which combines details from Rene Guenon, Julius Evola, and Plato.\footnote{The English translation of \textit{Revolt Against the Modern World} suggests the body parts for workers are the sexual and execratory organs. The Italian text mentions ``sesso e nutrimento'', not ``escrezione''. Plato refers to the ``belly''. Let that serve as a warning to those who draw conclusions based solely on translation.}

\begin{table}[h]\footnotesize\centering
\begin{tabular}{*{7}{p{.1\textwidth}}}\toprule
Caste & Sanskrit Name & Soul Power & Body Part & Vedas & Goal & Twice Born \\\midrule

Spiritual caste & Brahman & Spirit, nous & Head & Mouth & Transcendental Man & Yes \\

Nobility \& warriors & Kshatriya & Soul, thumos & Chest & Arms & True Man & Yes \\

Mercantile class \& Skilled workers & Vaishya & Desire, appetite & Belly & Thighs & Active Transformation of Material World & Yes \\

Servile workers & Shudra & Mechanical & ~ & Under foot & Passive transformation of Material World & No \\\bottomrule
\end{tabular}
\label{tab:CharacteristicsCaste}
\caption{Characteristics of Castes}
\end{table}

First of all, let us clear up some common misconceptions that are unfortunately often repeated and believed. The castes do not refer to race, since, within a racial group, there will be members of all four castes. Nor is caste membership inherited; although in practice, it may tend to that, it is not at all necessary in principle. Finally, gender is not part of castes; specifically, the notion that the brahman caste is inherently feminine and the kshatriya caste is masculine is false.

Caste membership refers only to one's function in society, at least in a Traditional society, since in the modern world, the caste memberships are all confused. Of course, one's societal function is a consequence of one's transcendental nature. Thus, caste is metaphysically necessary and not the result of allegedly oppressive political structures subject to the wills of men. When people perform their functions skillfully and willingly with understanding of the principle of castes, the society is stable. When the people become ignorant of that doctrine, resentment grows and class conflict is the result.

The first three castes have the possibility to be twice-born, and have their respective orders: spiritual orders, chivalric orders, and craft orders or guilds. Due to the differences in their respective natures, the nature of the orders will differ.

Evola mistakenly detects some incoherence in the social organization of the Middle Ages when he notes the different perspectives and moralities of the castes, when this is perfectly normal as our quote from William Blake points out. Even the pagan Plato in the \emph{Republic} proposed different rules of behavior for the guardians, warriors and everyone else. This is why any discussion about the tenets of a religion has to include the caste which accepts them. For example, the spiritual caste will know metaphysical doctrines in a trans-formal (image-less) way. The lower castes will instead believe the doctrines, which need to be presented using sensory (formal) imagery.

The roles and spiritual attitudes of the individual castes will be presented later.


\flrightit{Posted on 2011-02-03 by Cologero}

\begin{center}* * *\end{center}

\begin{footnotesize}
\begin{sffamily}
\texttt{Matt on 2011-02-04 at 15:53 said:}

``Let that serve as a warning to those who draw conclusions based solely on translation.''

Yes, I wish there were better english translations of Evola's and others' works. Unfortunate.

\hfill

\texttt{Perennial on 2011-02-05 at 02:34 said:}

Boy, isn't that the truth. Some are barely readable.

\hfill

\texttt{VisionsOfGlory14 on 2011-02-05 at 10:18 said:}

I don't know if it's bad translation or Evola's style but it takes some training to read him. He has an elevated language though that is consistent with his elevated view.

\hfill

\texttt{Cologero on 2011-02-05 at 11:56 said:}

His style is indeed difficult at first. He writes Italian as though it were Latin, using complex sentences and even obsolete words. Native Italians speakers have difficulty in reading him. (Difficulty? I really mean that can't even understand individual sentences.) That is why his books intended for a popular audience (Men among the ruins, Riding the tiger) are more widely read. However, they are hardly comprehensible without keeping Evola's entire worldview in mind. The difficulty in translating Evola is trying to maintain his style while maintaining readability. Furthermore, the translator has to understand metaphysics and tradition, more so, even, than the native language.

It's unfortunate, but it is hardly worth the trouble to translated his books. Financially, it makes no sense. Perhaps a professor could do it, but Evola is hardly popular in the Academy. I know firsthand because a visiting professor of a local university from Sicily publicly chewed me out. Oddly, everything she considered a fault, I considered a virtue. There is no grant money available. at least as far as my own inquiries have gone.

His three large volumes on philosophy take German Idealism to their ultimate conclusion in metaphysics, properly speaking. But they won't be translated. And I don't believe any student has done his PhD thesis on those books. His books on race and politics are interesting for their historical value, if for no other reason. I don't see any of these books being translated. There are few in the world even capable of doing so, and for us, it is hardly worth the effort.

\hfill

\texttt{Will on 2011-02-05 at 20:27 said:}

The folks over at Arktos, who published translations of Metaphysics of War and The Path of Cinnabar, have indicated on their website that more translations of Evola are forthcoming from them. They asked people via Facebook what other works people would like to see published, and the responses I saw were overwhelmingly for more Evola. So it would seem that there is at least some demand, though probably not much money to be made from it.

In regards to style, I can't comment on the accuracy of the various translators since I don't speak Italian, but in terms of readability, I've always thought that Evola reads easier than any of the major Traditionalist writers. His prose is more lively and forceful than Guenon or Coomarswamy or Schuon, all of whom seem rather dry in comparison to Evola, though no less profound.

\hfill

\texttt{JA on 2013-08-02 at 15:50 said:}

As I long term goal, I would actually like to write my PhD thesis on Evola's contributions to aesthetics and art theory.


\hfill

\end{sffamily}
\end{footnotesize}

